% !TEX program = xelatex
\documentclass[10pt,letterpaper,twocolumn]{article}

\usepackage[top=0.85in, bottom=0.85in, left=0.75in, right=0.75in]{geometry}
\usepackage{fontspec}
\usepackage{unicode-math}
\setmainfont{Latin Modern Roman}[
  Extension=.otf,
  UprightFont=lmroman10-regular,
  BoldFont=lmroman10-bold,
  ItalicFont=lmroman10-italic,
  BoldItalicFont=lmroman10-bolditalic,
]
\setsansfont{Latin Modern Sans}[
  Extension=.otf,
  UprightFont=lmsans10-regular,
  BoldFont=lmsans10-bold,
  ItalicFont=lmsans10-oblique,
  BoldItalicFont=lmsans10-boldoblique,
]
\newfontfamily\acbold{ywft-processing-bold}[Path=../../system/public/type/webfonts/,Extension=.ttf]
\newfontfamily\aclight{ywft-processing-light}[Path=../../system/public/type/webfonts/,Extension=.ttf]
\setmonofont{Latin Modern Mono}[
  Extension=.otf,
  UprightFont=lmmono10-regular,
  ItalicFont=lmmono10-italic,
  Scale=0.85,
]

\usepackage{xcolor}
\usepackage{titlesec}
\usepackage{enumitem}
\usepackage{booktabs}
\usepackage{fancyhdr}
\usepackage{graphicx}
\graphicspath{{figures/}{../../papers/arxiv-ac/figures/}}
\usepackage{ragged2e}
\usepackage{microtype}
\usepackage{eso-pic}
\usepackage{tikz}
\usepackage{hyperref}
\usepackage{xurl}

\definecolor{acpink}{RGB}{180,72,135}
\definecolor{acpurple}{RGB}{120,80,180}
\definecolor{acdark}{RGB}{64,56,74}
\definecolor{acgray}{RGB}{119,119,119}
\definecolor{palspink}{RGB}{205,92,155} % #cd5c9b — canonical purple-pals.svg fill

% === PALS SIDE-STAMPS (mirrors bios.mjs side-stamp pattern + ac-paper-layout.sty) ===
\AddToShipoutPictureBG{%
  \begin{tikzpicture}[remember picture, overlay]
    \pgfmathsetmacro{\palswiggle}{7*sin(\thepage * 67)}%
    \node[opacity=0.10, anchor=center, rotate={90 + \palswiggle}]
      at ([xshift=-1.0cm, yshift=-0.33\paperheight]current page.north west)
      {
\includegraphics[width=13cm]{palsmark}};
    \node[opacity=0.10, anchor=center, rotate={-90 + \palswiggle}]
      at ([xshift=1.0cm, yshift=-0.33\paperheight]current page.north east)
      {
\includegraphics[width=13cm]{palsmark}};
  \end{tikzpicture}%
}

\hypersetup{colorlinks=true,linkcolor=acpurple,urlcolor=acpurple,citecolor=acpurple,
  pdftitle={The Pals Mark: Una historia del logo de Aesthetic.Computer}}

\titleformat{\section}{\normalfont\bfseries\normalsize\uppercase}{\thesection.}{0.5em}{}
\titlespacing{\section}{0pt}{1.2em}{0.3em}
\titleformat{\subsection}{\normalfont\bfseries\small}{\thesubsection}{0.5em}{}
\titlespacing{\subsection}{0pt}{0.8em}{0.2em}

\pagestyle{fancy}\fancyhf{}
\renewcommand{\headrulewidth}{0pt}
\fancyhead[C]{\footnotesize\color{acpink}\textit{The Pals Mark --- \ac{}}}
\fancyfoot[C]{\footnotesize\thepage}

\setlist[itemize]{nosep, leftmargin=1.2em, itemsep=0.15em}
\setlist[enumerate]{nosep, leftmargin=1.4em}
\setlength{\columnsep}{1.8em}
\setlength{\parindent}{1em}
\setlength{\parskip}{0.3em}

\newcommand{\ac}{\textsc{Aesthetic}{\color{acpink}.}\textsc{Computer}}
\newcommand{\acdot}{{\color{acpink}.}}
\newcommand{\pals}{\textsc{Pals}}
% commit link: \commit{shorthash}{message} -> links the hash to GitHub
\newcommand{\commit}[2]{\href{https://github.com/whistlegraph/aesthetic-computer/commit/#1}{\texttt{#1}}, \emph{``#2''}}
% repo file link
\newcommand{\repofile}[1]{\href{https://github.com/whistlegraph/aesthetic-computer/blob/main/#1}{\texttt{#1}}}

% Cover language bar. The same paper is published in six languages; each
% edition links to the others by canonical permalink. \klanghere marks the
% current language; \klang links a sibling edition by its URL suffix.
\newcommand{\klang}[2]{\href{https://papers.aesthetic.computer/pals-mark-26-arxiv#2.pdf}{\color{acgray}#1}}
\newcommand{\klanghere}[1]{{\bfseries\color{acpink}#1}}
\newcommand{\klangsep}{{\color{acgray!60}\,\textperiodcentered\,}}

\begin{document}

\twocolumn[{%
\begin{center}
{\acbold\fontsize{27pt}{31pt}\selectfont\color{acdark} The \pals{} Mark}\par
\vspace{0.5em}
\includegraphics[width=0.6\textwidth]{figures/cover}\par
\vspace{0.4em}
{\itshape\fontsize{14.5pt}{17pt}\selectfont\color{acpink} Una historia del logo de Aesthetic{\acdot}Computer}\par
\vspace{0.5em}
{\small\color{acgray} \href{https://prompt.ac/@jeffrey}{@jeffrey} \,$\cdot$\, Aesthetic{\acdot}Computer \,$\cdot$\, logo de Niki Stebbins \,$\cdot$\, ORCID \href{https://orcid.org/0009-0007-4460-4913}{0009-0007-4460-4913}}\par
\vspace{0.45em}
{\footnotesize\sffamily\klang{EN}{}\klangsep\klang{DA}{-da}\klangsep\klanghere{ES}\klangsep\klang{ZH}{-zh}\klangsep\klang{JA}{-ja}\klangsep\klang{RU}{-ru}}\par
\vspace{0.55em}
{\color{acpink}\rule{\textwidth}{1.2pt}}
\vspace{0.5em}
\end{center}

\begin{quote}
\small\noindent\textbf{Resumen.}
Los \pals{} son el logo de \ac{}: dos figuras dibujadas a mano que se toman
de la mano con los brazos en alto. Este artículo relata cómo se creó la
marca --- encargada a la artista Niki Stebbins y dibujada, durante el verano
de 2023, \emph{dentro} de \ac{} usando su propia herramienta de whistlegraph
\texttt{wgr} --- y cómo luego se convirtió en infraestructura: un
``capturador de logos'' rotatorio en \texttt{pals.aesthetic.computer} que
sirve las muchas poses de la marca a través de avatares, artículos, sistema
operativo nativo y la web, cada una un enlace de vuelta al proyecto.
\end{quote}
\vspace{0.5em}
}]

\section{Qué es \ac{}}

\ac{} es un entorno de ejecución y una red social pensados primero para
móvil, orientados a la computación creativa --- una única barra de URL que
funciona además como una interfaz al estilo de un instrumento musical, donde
comandos memorizables invocan programas interactivos llamados ``piezas''.
Admite tanto JavaScript (\texttt{.mjs}) como un dialecto propio de Lisp,
\textbf{KidLisp} (118 funciones integradas). Todo el sistema se comunica con
las piezas a través de una única API central (\texttt{disk.mjs},
$\sim$600\,KB).

\textbf{Cuándo se puso en línea.} El dominio \texttt{aesthetic.computer} se
mantiene desde el \textbf{2013-12-12}. La plataforma tal como existe fue
concebida a partir de \emph{No Paint} (2020), el desarrollo comenzó en
\textbf{2021}, y se ha desplegado de forma continua desde el repositorio
actual desde \textbf{diciembre de 2022} ($\sim$20{,}300 commits).

\textbf{Estadísticas de la plataforma} (en vivo, \today):
\begin{itemize}
  \item $\sim$400 piezas integradas en el repositorio (383 \texttt{.mjs} + 20
        \texttt{.lisp}); 265 obras publicadas por usuarios.
  \item \textbf{17{,}538} programas de KidLisp almacenados en producción ---
        un corpus creciente de arte generativo creado por la comunidad.
  \item Federación ATProto / Bluesky; colecciones de registros por handle
        (pinturas, estados de ánimo, piezas, KidLisp) en la red
        \texttt{at://}.
\end{itemize}

\textbf{A quién sirve, y quién conoce la marca.} Los usuarios de \ac{} son una
comunidad de programadores creativos, artistas y ``Pals'' no técnicos que
publican al instante piezas bajo \texttt{@handle}. La plataforma atrae a un
amplio público en todo el campo del arte y la tecnología, y obras de \ac{}
han entrado en colecciones institucionales. La propia marca \pals{} recorre
cada una de estas superficies: la pantalla de arranque, los artículos, los
avatares sociales, el sistema operativo nativo. Su linaje conceptual en la
comunidad de música radical por ordenador de \emph{Goodiepal \& Pals} (véase
la sección de linaje más abajo) da al nombre una historia reconocida en el
arte de los medios más allá del propio público de \ac{}.

\section{Qué representa la marca}

El logo \pals{} es \textbf{dos figuras humanas de contorno que se toman de la
mano con los brazos en alto} --- como dos amigos que caminan o bailan juntos
--- trazadas con una línea suelta, orgánica, dibujada a mano. El vector
estático canónico (\texttt{system/public/purple-pals.svg}) es de color
rosa/magenta, \texttt{\#cd5c9b}. \emph{No} es ``un par de ojos'', una lectura
errónea anterior que desde entonces se ha corregido en el conjunto de
artículos.

De manera crucial, la marca se hizo \emph{con} la plataforma, no meramente
\emph{para} ella: Niki dibujó las exploraciones del logo usando \texttt{wgr},
la pieza de dibujo renderizadora de whistlegraphs que corre dentro del propio
\ac{}. El logo es, por tanto, a la vez una identidad y una demostración de la
herramienta.

\section{Cómo se creó la marca}

\subsection{Primero, un sello}
Comienza en el verano de \textbf{2020} con un pequeño regalo. Niki Stebbins
envía por correo un dibujo --- una pequeña estrella púrpura --- y Jeffrey lo
incorpora al juego como un sello bajo el handle \texttt{@niki} y pide más.
Durante los dos años siguientes la colaboración es una \emph{sesión de dibujo}
suelta y continua: sellos, un juego de dibujo, garabatos de whistlegraph.
Para \textbf{2022} ya es trabajo de diseño remunerado. Nada de ello es aún el
logo, pero es el suelo del que el logo crece --- el hábito de dos personas
dibujando juntas dentro de la herramienta.

\subsection{Reabrir el diseño (junio de 2023)}
El logo propiamente dicho empieza con una insatisfacción. En \textbf{junio de
2023} Niki escribe que no está ``100\% satisfecha con el logo que te dejé'', y
lo reabre. Piensa en voz alta sobre el color --- ``como púrpura o
verde\dots{} azul no, hay demasiados logos de apps azules'' --- y sobre la
forma: ¿debería ser la marca ``como un marco, solo una forma cerrada'', algo
dentro de lo cual pudieras trabajar? Envía una página de bocetos con
candidatos y elige un favorito --- ``el del círculo es el que más me gusta''
--- y de inmediato le preocupa que sea ``bastante parecido al logo de
are.na'', y sigue adelante. La instrucción clave que se da a sí misma es una
corrección de proporción: empujarlo ``para que parezca menos estrellas y más
figuras''. Una página de los estudios recibe su propio veredicto --- ``parecen
estrellas de mar jaja'' --- y para arreglarlo abandona la pantalla por
completo, atando y fotografiando tramos de \textbf{cuerda} en un parque para
entender cómo caen de verdad dos cuerpos enlazados.

\subsection{Construir las dos figuras (julio--agosto de 2023)}
A lo largo del verano la marca se resuelve en \emph{dos} figuras en lugar de
una. Niki compra una cámara digital \textbf{Kodak EasyShare} de segunda mano
específicamente para fotografiar el logo, ajusta una versión al ``tamaño
exacto del logo del VS Marketplace'', y pregunta si Jeffrey quiere ``más
dibujos del logo en a.c para fotos de perfil en redes sociales y la cosa de la
app-store''. El repositorio concuerda con los mensajes: en
\commit{5bcefc5a9}{better defaults in \texttt{wgr} for niki's logo efforts}
la propia herramienta de dibujo se afina para ella, y en
\commit{a0009f5d0}{add new vscode extension logo from niki} aterriza un
recurso ya distribuido. A finales de agosto informa de que están terminadas
las ``variaciones de los \textbf{dos logos}'', describiendo el gesto con
precisión --- ``la figura de la derecha tiraba de la de la izquierda\dots{} el
brazo de la figura izquierda está más alto''. La marca es ahora lo que sigue
siendo: dos personas, con los brazos en alto, una tirando de la otra. El
\textbf{2023-08-31} se elige una foto de perfil de Instagram.

\subsection{Helado, y la puesta en marcha (septiembre de 2023)}
El final es la parte más extraña y la mejor. Para fotografiar la marca en el
mundo, Niki la dibuja a tamaño real sobre el pavimento e intenta capturarla
con helado --- pero el helado del estadio viene en vasitos, así que se desvía
a un McDonald's por un cono y consigue la foto. El \textbf{2023-09-09} ``envió
los logos al correo de a.c'' con el asunto \emph{!c3 cr34m l0go\$} (cuerpo:
``ice cream logos 8)''). Esa misma semana la marca deja de ser una imagen y se
convierte en infraestructura: aterriza \commit{0aba1a378}{add dynamic logo
generator}, y Jeffrey informa del resultado en una línea --- ``\^{}\ devolverá
ahora un logo aleatorio cada vez''. El logo se dibujó \emph{con} la plataforma
(en la herramienta de whistlegraph \texttt{wgr}) y luego fue servido \emph{por}
ella. Nunca fue solo una imagen; siempre fue también un programa.

\section{El endpoint: \texttt{pals.aesthetic.computer}}

La marca se convirtió en infraestructura en septiembre de 2023, cuando \ac{}
levantó un servicio de logos dedicado y le dio un hogar memorizable.

\begin{itemize}
  \item \textbf{2023-09-05} --- \commit{0aba1a378}{add dynamic logo
        generator}. El módulo de backend \repofile{system/backend/logo.mjs}
        lleva la marca de tiempo \texttt{23.09.05.22.31}.
  \item \textbf{2023-09-12} --- \commit{f1a4f272c}{logo -> pals.aesthetic}
        --- la ruta se renombra a la identidad \pals{}.
  \item \textbf{2023-09-13} --- \commit{12e348035}{pals.aesthetic.computer
        url fix} --- el endpoint está en vivo y es accesible en su propio
        subdominio.
\end{itemize}

\noindent Así que \href{https://pals.aesthetic.computer}{\texttt{pals.aesthetic.computer}}
\textbf{se puso en línea el 2023-09-12--13}, tres días después de que
aterrizara el generador dinámico.

\subsection{Cómo funciona como ``capturador de logos''}

El endpoint es un diminuto proxy rotatorio, no una única imagen fija. Cada
solicitud devuelve una pose \pals{} distinta:

\begin{enumerate}
  \item \href{https://pals.aesthetic.computer}{\texttt{pals.aesthetic.computer}}
        redirige (mediante \repofile{system/netlify.toml}) a
        \href{https://aesthetic.computer/api/logo}{\texttt{aesthetic.computer/api/logo}}.
  \item \repofile{system/backend/logo.mjs} mantiene una lista de
        \emph{pinturas} \pals{} en el bucket CDN
        \texttt{pals-aesthetic-computer} y devuelve una al azar en cada
        llamada (evitando una repetición inmediata).
  \item La función negocia el contenido: un navegador recibe una página HTML
        a sangre completa sobre un púrpura aleatorio, que se puede tocar para
        pasar a la siguiente pose (pasando \texttt{?previousLogo=} para que
        nunca se repita); una solicitud de
        \texttt{curl}/\texttt{<img>}/servidor recibe directamente el
        \texttt{.png} en bruto. La misma URL es así a la vez un visor y un
        recurso enlazable directamente.
\end{enumerate}

\subsection{Las muchas poses de los \pals{}}

El conjunto de poses del capturador es una galería de pinturas \texttt{wgr} de
mediados de 2023 (la serie \texttt{painting-2023.7.29}--\texttt{2023.8.21}).
Como el logo es una \emph{colección de poses dibujadas} en lugar de un único
glifo congelado, la marca se percibe como viva: cada carga de página,
actualización de avatar o petición de favicon puede mostrar a las dos figuras
en una postura ligeramente distinta. El conjunto completo de doce poses en
vivo se reproduce en el Apéndice~A.

\section{Uso en distintas propiedades}

Como cualquier propiedad puede enlazar directamente
\href{https://pals.aesthetic.computer}{\texttt{pals.aesthetic.computer}}
(o incrustar el estático
\repofile{system/public/purple-pals.svg}), una única marca hilvana todo el
ecosistema, y cada instancia es un enlace de vuelta al proyecto:

\begin{itemize}
  \item \textbf{Avatares sociales y Gravatar} --- el endpoint rotatorio
        suministra una imagen de perfil que cambia de pose de forma discreta,
        de modo que la misma identidad aparece en distintos servicios mientras
        siempre resuelve a \ac{}.
  \item \textbf{Web y marketing} --- favicons, encabezados de documentación y
        etiquetas \texttt{<img>} incrustadas extraen los \pals{} en vivo; un
        toque cambia de pose, convirtiendo el propio logo en un pequeño
        anticipo interactivo de la plataforma.
  \item \textbf{Artículos} --- los \pals{} recorren los márgenes de todo
        borrador de trabajo de \ac{} como sellos laterales semiopacos que se
        contonean página a página (véase esta página), y aparecen en los
        colofones de las ediciones impresas.
  \item \textbf{Sistema operativo nativo y escritorio} --- los \pals{} sirven
        como identidad de arranque (splash de Plymouth) en AC Native OS y como
        el icono vectorial de la bandeja en AC Electron y MenuBand ---
        teñibles con el color de acento y con posibilidad de insignia.
  \item \textbf{Impresión (pendiente)} --- en mayo de 2025 Jeffrey contactó
        con Casey Rubber Stamps para tallar sellos físicos de los \pals{} en
        tres tamaños; el pedido se estancó y nunca se completó.
\end{itemize}

\section{¿Por qué ``\pals{}''? El linaje Goodiepal}

El nombre no es decoración; es una afirmación estructural, y su peso viene de
fuera de \ac{}. Como se argumenta en el artículo de \ac{}
\emph{Radical Computer Art: Goodiepalian Approaches in Aesthetic
Computer} (\repofile{papers/arxiv-goodiepal/goodiepal.tex}) y en el artículo
de la plataforma que este cita, ``\pals{}'' desciende de \textbf{Goodiepal \&
Pals} --- el laxo grupo de afinidad en torno al músico danés-feroés y teórico
de la Música Radical por Ordenador Goodiepal (Gaeoudji SYGNOK, n.~1974).

El modelo \pals{} es una política de la organización: un grupo sin membresía,
sin respaldo institucional y sin estructura formal, mantenido unido por la
práctica compartida y el reconocimiento mutuo antes que por la afiliación.
Goodiepal enseñó Historia y Estética de la Música Electrónica en la Real
Academia Danesa de Música y la dejó en 2008 tras un conflicto descrito como
tener ``el tipo equivocado de amigos'', habiendo declarado la ``guerra
intelectual contra la estupidez en la música por ordenador y el arte de los
medios modernos''. Jeffrey hizo una gira con Goodiepal \& Pals por Europa en
2018, y las herramientas construidas para esas actuaciones son ancestros
directos de la arquitectura de \ac{}.

Así, el logo lleva esa historia política del arte de los medios a la marca.
Como lo plantea el artículo de Goodiepal, nombrar el logo ``\pals{}'' --- y
titular la partitura \emph{Score for Aesthetic{\color{acpink}.}Computer \&
Pals} --- ``no es branding; es una afirmación estructural: la plataforma
pertenece a su comunidad de practicantes, no a su desarrollador ni a sus
inversores''. Dos figuras que se toman de la mano, hechas \emph{dentro} de la
herramienta, es el argumento representado como marca: \ac{} es \emph{para}
cualquiera que se presente --- para los \pals{}.

\section{Autoría}

Jeffrey contrató a \textbf{Niki Stebbins} para diseñar el logo de \ac{} bajo
su dirección. A lo largo de 2022--2023 ella desarrolló la marca dentro de la
propia herramienta de dibujo \texttt{wgr} de la plataforma, atravesando las
exploraciones del cuaderno de bocetos, la cuerda y el helado documentadas más
arriba y entregando el dispositivo de dos figuras que se distribuye como
\texttt{purple-pals.svg}. El capturador rotatorio
\texttt{pals.aesthetic.computer} se puso en vivo el 2023-09-12--13; el vector
estático canónico aparece por primera vez el 2024-01-31
(\commit{bf29eadf8}{host purple-pals graphic for docs}); y a partir de 2026-02
el repositorio y los artículos sitúan explícitamente a ``\pals{}'' en el
centro como identidad de la plataforma (\commit{48dd8923e}{rename score to
``Score for Aesthetic.Computer \& Pals''}). La marca es una obra por encargo:
creada por Niki, dirigida por Jeffrey, y propiedad de \ac{}.

\vspace{0.6em}
{\footnotesize\color{acgray}\noindent Fuentes: historial de commits del
repositorio;
\repofile{system/backend/logo.mjs};
\repofile{system/netlify/functions/logo.js};
\repofile{system/netlify.toml};
\repofile{papers/arxiv-goodiepal/goodiepal.tex};
\repofile{system/public/purple-pals.svg}.\par}

% ============================ APPENDICES ============================
\onecolumn
\appendix

\noindent\emph{Nota sobre las fechas.} Cada recurso visual a continuación
lleva como pie la fecha en que se creó: las poses de \texttt{wgr} por la marca
de tiempo codificada en el slug de la pintura; los recursos del monorepo por
su fecha de primer commit (git \texttt{--diff-filter=A}); las capturas en vivo
por la fecha en que se fotografiaron desde producción.

\section*{Apéndice A --- Las muchas poses (conjunto en vivo de \texttt{pals.aesthetic.computer})}
\addcontentsline{toc}{section}{Apéndice A --- Poses}
\noindent Las doce pinturas \texttt{wgr} por las que rota el capturador de
logos, extraídas en vivo del CDN \texttt{pals-aesthetic-computer}. El mismo
motivo de dos figuras, cada una con una postura y un tratamiento de color
distintos.
\vspace{1.2em}

\begin{minipage}[t]{0.242\textwidth}\centering
  
\includegraphics[width=\linewidth,height=\linewidth,keepaspectratio]{figures/poses/small/pose-1.png}\\[2pt]
  {\scriptsize\textbf{pose 1}}\\[-1pt]
  {\scriptsize\color{acgray}2023-08-01}
\end{minipage}\hfill\begin{minipage}[t]{0.242\textwidth}\centering
  \includegraphics[width=\linewidth,height=\linewidth,keepaspectratio]{figures/poses/small/pose-2.png}\\[2pt]
  {\scriptsize\textbf{pose 2}}\\[-1pt]
  {\scriptsize\color{acgray}2023-08-21}
\end{minipage}\hfill\begin{minipage}[t]{0.242\textwidth}\centering
  
\includegraphics[width=\linewidth,height=\linewidth,keepaspectratio]{figures/poses/small/pose-3.png}\\[2pt]
  {\scriptsize\textbf{pose 3}}\\[-1pt]
  {\scriptsize\color{acgray}2023-08-21}
\end{minipage}\hfill\begin{minipage}[t]{0.242\textwidth}\centering
  \includegraphics[width=\linewidth,height=\linewidth,keepaspectratio]{figures/poses/small/pose-4.png}\\[2pt]
  {\scriptsize\textbf{pose 4}}\\[-1pt]
  {\scriptsize\color{acgray}2023-08-01}
\end{minipage}

\vspace{0.8em}

\begin{minipage}[t]{0.242\textwidth}\centering
  \includegraphics[width=\linewidth,height=\linewidth,keepaspectratio]{figures/poses/small/pose-5.png}\\[2pt]
  {\scriptsize\textbf{pose 5}}\\[-1pt]
  {\scriptsize\color{acgray}2023-08-21}
\end{minipage}\hfill\begin{minipage}[t]{0.242\textwidth}\centering
  \includegraphics[width=\linewidth,height=\linewidth,keepaspectratio]{figures/poses/small/pose-6.png}\\[2pt]
  {\scriptsize\textbf{pose 6}}\\[-1pt]
  {\scriptsize\color{acgray}2023-08-01}
\end{minipage}\hfill\begin{minipage}[t]{0.242\textwidth}\centering
  \includegraphics[width=\linewidth,height=\linewidth,keepaspectratio]{figures/poses/small/pose-7.png}\\[2pt]
  {\scriptsize\textbf{pose 7}}\\[-1pt]
  {\scriptsize\color{acgray}2023-08-21}
\end{minipage}\hfill\begin{minipage}[t]{0.242\textwidth}\centering
  \includegraphics[width=\linewidth,height=\linewidth,keepaspectratio]{figures/poses/small/pose-8.png}\\[2pt]
  {\scriptsize\textbf{pose 8}}\\[-1pt]
  {\scriptsize\color{acgray}2023-08-21}
\end{minipage}

\vspace{0.8em}

\begin{minipage}[t]{0.242\textwidth}\centering
  \includegraphics[width=\linewidth,height=\linewidth,keepaspectratio]{figures/poses/small/pose-9.png}\\[2pt]
  {\scriptsize\textbf{pose 9}}\\[-1pt]
  {\scriptsize\color{acgray}2023-08-01}
\end{minipage}\hfill\begin{minipage}[t]{0.242\textwidth}\centering
  \includegraphics[width=\linewidth,height=\linewidth,keepaspectratio]{figures/poses/small/pose-10.png}\\[2pt]
  {\scriptsize\textbf{pose 10}}\\[-1pt]
  {\scriptsize\color{acgray}2023-08-01}
\end{minipage}\hfill\begin{minipage}[t]{0.242\textwidth}\centering
  \includegraphics[width=\linewidth,height=\linewidth,keepaspectratio]{figures/poses/small/pose-11.png}\\[2pt]
  {\scriptsize\textbf{pose 11}}\\[-1pt]
  {\scriptsize\color{acgray}2023-07-31}
\end{minipage}\hfill\begin{minipage}[t]{0.242\textwidth}\centering
  
\includegraphics[width=\linewidth,height=\linewidth,keepaspectratio]{figures/poses/small/pose-12.png}\\[2pt]
  {\scriptsize\textbf{pose 12}}\\[-1pt]
  {\scriptsize\color{acgray}2023-07-29}
\end{minipage}

\vspace{0.8em}




\newpage
\section*{Apéndice B --- La marca a través del monorepo}
\addcontentsline{toc}{section}{Apéndice B --- Recursos del monorepo}
\noindent Los mismos \pals{} dibujados una vez y revestidos de nuevo en todas
partes: el vector canónico y sus descendientes por toda la base de código ---
extensión del editor, splash de arranque del sistema operativo nativo, bandeja
de escritorio, portadas de pop, sello de facturas y el conjunto de
combinaciones de color de keymaps. Cada mosaico lleva como pie su fecha de
primer commit.
\vspace{1.2em}

\begin{minipage}[t]{0.194\textwidth}\centering
  
\includegraphics[width=\linewidth,height=\linewidth,keepaspectratio]{figures/assets/c01-purple-pals.png}\\[2pt]
  {\scriptsize\textbf{purple-pals.svg}}\\[-1pt]
  {\scriptsize\color{acgray}2024-01-31}
\end{minipage}\hfill\begin{minipage}[t]{0.194\textwidth}\centering
  
\includegraphics[width=\linewidth,height=\linewidth,keepaspectratio]{figures/assets/c02-vscode.png}\\[2pt]
  {\scriptsize\textbf{vscode-extension}}\\[-1pt]
  {\scriptsize\color{acgray}2024-01-31}
\end{minipage}\hfill\begin{minipage}[t]{0.194\textwidth}\centering
  
\includegraphics[width=\linewidth,height=\linewidth,keepaspectratio]{figures/assets/c03-tvos.png}\\[2pt]
  {\scriptsize\textbf{apple/tvos-tapes}}\\[-1pt]
  {\scriptsize\color{acgray}2024-01-31}
\end{minipage}\hfill\begin{minipage}[t]{0.194\textwidth}\centering
  
\includegraphics[width=\linewidth,height=\linewidth,keepaspectratio]{figures/assets/c04-plymouth.png}\\[2pt]
  {\scriptsize\textbf{AC-OS boot splash}}\\[-1pt]
  {\scriptsize\color{acgray}2026-02-21}
\end{minipage}\hfill\begin{minipage}[t]{0.194\textwidth}\centering
  
\includegraphics[width=\linewidth,height=\linewidth,keepaspectratio]{figures/assets/c05-tray.png}\\[2pt]
  {\scriptsize\textbf{Electron tray}}\\[-1pt]
  {\scriptsize\color{acgray}2026-05-20}
\end{minipage}

\vspace{0.8em}

\begin{minipage}[t]{0.194\textwidth}\centering
  
\includegraphics[width=\linewidth,height=\linewidth,keepaspectratio]{figures/assets/c14-stamp.png}\\[2pt]
  {\scriptsize\textbf{invoice stamp}}\\[-1pt]
  {\scriptsize\color{acgray}2026-03-06}
\end{minipage}\hfill\begin{minipage}[t]{0.194\textwidth}\centering
  
\includegraphics[width=\linewidth,height=\linewidth,keepaspectratio]{figures/assets/c12-gray.png}\\[2pt]
  {\scriptsize\textbf{grayscale (quotes)}}\\[-1pt]
  {\scriptsize\color{acgray}2026-05-21}
\end{minipage}\hfill\begin{minipage}[t]{0.194\textwidth}\centering
  
\includegraphics[width=\linewidth,height=\linewidth,keepaspectratio]{figures/assets/c13-hellsine.png}\\[2pt]
  {\scriptsize\textbf{pop cover glyph}}\\[-1pt]
  {\scriptsize\color{acgray}2026-05-24}
\end{minipage}\hfill\begin{minipage}[t]{0.194\textwidth}\centering
  
\includegraphics[width=\linewidth,height=\linewidth,keepaspectratio]{figures/assets/c06-pink.png}\\[2pt]
  {\scriptsize\textbf{keymaps pink}}\\[-1pt]
  {\scriptsize\color{acgray}2026-06-13}
\end{minipage}\hfill\begin{minipage}[t]{0.194\textwidth}\centering
  
\includegraphics[width=\linewidth,height=\linewidth,keepaspectratio]{figures/assets/c07-blue.png}\\[2pt]
  {\scriptsize\textbf{keymaps blue}}\\[-1pt]
  {\scriptsize\color{acgray}2026-06-13}
\end{minipage}

\vspace{0.8em}

\begin{minipage}[t]{0.194\textwidth}\centering
  
\includegraphics[width=\linewidth,height=\linewidth,keepaspectratio]{figures/assets/c08-green.png}\\[2pt]
  {\scriptsize\textbf{keymaps green}}\\[-1pt]
  {\scriptsize\color{acgray}2026-06-13}
\end{minipage}\hfill\begin{minipage}[t]{0.194\textwidth}\centering
  
\includegraphics[width=\linewidth,height=\linewidth,keepaspectratio]{figures/assets/c09-orange.png}\\[2pt]
  {\scriptsize\textbf{keymaps orange}}\\[-1pt]
  {\scriptsize\color{acgray}2026-06-13}
\end{minipage}\hfill\begin{minipage}[t]{0.194\textwidth}\centering
  
\includegraphics[width=\linewidth,height=\linewidth,keepaspectratio]{figures/assets/c10-purple.png}\\[2pt]
  {\scriptsize\textbf{keymaps purple}}\\[-1pt]
  {\scriptsize\color{acgray}2026-06-13}
\end{minipage}\hfill\begin{minipage}[t]{0.194\textwidth}\centering
  
\includegraphics[width=\linewidth,height=\linewidth,keepaspectratio]{figures/assets/c11-yellow.png}\\[2pt]
  {\scriptsize\textbf{keymaps yellow}}\\[-1pt]
  {\scriptsize\color{acgray}2026-06-13}
\end{minipage}\hfill

\newpage
\section*{Apéndice C --- Emplazamientos en vivo (interfaces reales)}
\addcontentsline{toc}{section}{Apéndice C --- Emplazamientos en vivo}
\noindent Los \pals{} fotografiados viviendo en interfaces de producción
reales, cada una con como pie su URL de origen (los enlaces están en vivo): la
ficha del Visual Studio Marketplace (los \pals{} como icono de producto de la
extensión), el README del monorepo tanto en GitHub como en el host canónico
Tangled, y el endpoint rotatorio
\href{https://pals.aesthetic.computer}{\texttt{pals.aesthetic.computer}}
(misma URL, distintas poses al recargar --- más el
\texttt{/api/logo.png} en bruto entregado a las etiquetas \texttt{<img>} y a
los avatares).
\vspace{1.0em}

% Hand-authored: the Pals living in real production UIs, each with a caption + live URL.

\newcommand{\uishot}[3]{%
  % #1 image  #2 caption  #3 url
  \begin{center}
  \fbox{\includegraphics[width=0.96\linewidth]{#1}}\\[0.4em]
  {\footnotesize #2}\\[0.1em]
  {\footnotesize\ttfamily\href{#3}{#3}}
  \end{center}
  \vspace{0.7em}
}
\newcommand{\uishottwo}[6]{%
  % (#1 img #2 cap #3 url) (#4 img #5 cap #6 url)
  \noindent
  \begin{minipage}[t]{0.485\linewidth}\centering
    \fbox{\includegraphics[width=0.97\linewidth]{#1}}\\[0.35em]
    {\footnotesize #2}\\[0.1em]{\scriptsize\ttfamily\href{#3}{#3}}
  \end{minipage}\hfill
  \begin{minipage}[t]{0.485\linewidth}\centering
    \fbox{\includegraphics[width=0.97\linewidth]{#4}}\\[0.35em]
    {\footnotesize #5}\\[0.1em]{\scriptsize\ttfamily\href{#6}{#6}}
  \end{minipage}
  \vspace{0.9em}
}

\uishot{figures/live/ui-vscode-marketplace.png}{%
  \textbf{Visual Studio Marketplace} --- the \pals{} as the product icon of
  the official \ac{} VS~Code extension (1{,}201 installs).}{%
  https://marketplace.visualstudio.com/items?itemName=aesthetic-computer.aesthetic-computer-code}

\uishottwo
  {figures/live/ui-github-readme.png}{\textbf{GitHub} --- the \pals{} atop the
   monorepo README.}{https://github.com/whistlegraph/aesthetic-computer}
  {figures/live/ui-tangled-readme.png}{\textbf{Tangled} (canonical git host)
   --- the same README, same mark.}{https://tangled.org/aesthetic.computer/core}

\uishot{figures/compare/pals-live-strip.png}{%
  \textbf{\texttt{pals.aesthetic.computer}} in a live browser --- the same URL
  reloaded shows different poses (the rotating ``logo grabber''); the raw
  \texttt{/api/logo.png} is what \texttt{<img>} tags and avatars (Gravatar,
  social) fetch across the web.}{%
  https://pals.aesthetic.computer}



\end{document}
